\subsection{Evaluation Methodology}
\label{sec:evaluation}

\subsubsection{Evaluation Goals}
\label{sec:eval-goals}

We evaluate Guard-Agent along three dimensions:

\begin{enumerate}
    \item \textbf{Correctness of checkpoint placement}: Can the approach correctly identify (a)~which data structures must be protected, (b)~where in the control flow to place checkpoint/restart calls, and (c)~under what conditions to trigger recovery?

    \item \textbf{Robustness across application diversity}: How does performance vary across applications with different data structure characteristics (fixed-size vs.\ variable-size state), computation patterns (iterative, task-parallel, pipeline, irregular), and codebase complexity (proxy apps to production frameworks)?

    \item \textbf{Efficiency}: How many LLM iterations, tokens, and wall-clock time are needed to produce a correct resilient version? How does guided analysis (Guard-Agent) compare to unassisted LLM modification (baseline)?
\end{enumerate}

%% ---------------------------------------------------------------------------
\subsubsection{Application Test-Bed}
\label{sec:testbed}

We construct a test-bed of 20 HPC applications organized in a two-dimensional classification matrix crossing \emph{computation pattern} (4~categories) with \emph{state characteristics} (2~categories), yielding 8~cells with 2--3 applications each.

\paragraph{Classification Matrix.}

Table~\ref{tab:classification} presents the full classification.

\begin{table}[t]
\centering
\caption{Application classification matrix: computation pattern $\times$ state characteristics.}
\label{tab:classification}
\small
\begin{tabular}{@{}lll@{}}
\toprule
\textbf{Computation Pattern} & \textbf{Fixed-size State} & \textbf{Variable-size State} \\
\midrule
Iterative / Time-stepping & CoMD~\cite{comd}, LULESH~\cite{lulesh}, miniFE~\cite{minife} & LAMMPS~\cite{lammps}, WarpX~\cite{warpx}, Nyx~\cite{nyx} \\
Task-parallel / DAG       & Chameleon~\cite{chameleon}, DPLASMA~\cite{dplasma}   & CLAMR~\cite{clamr}, deal.II~\cite{dealii} \\
Pipeline / Staged          & Palabos~\cite{palabos}, Nektar++~\cite{nektar}, SU2~\cite{su2} & RAxML-NG~\cite{raxml-ng}, ExaML~\cite{examl} \\
Irregular / Adaptive       & AMG~\cite{amg}, miniVite~\cite{minivite}        & SAMRAI~\cite{samrai}, Uintah~\cite{uintah}, AMReX~\cite{amrex} \\
\bottomrule
\end{tabular}
\end{table}

\paragraph{Iterative / Time-stepping.}
These applications advance state through a main loop with a fixed or adaptive time step. Checkpoint placement is typically at the loop boundary. \emph{Fixed-size} variants (CoMD~\cite{comd}, LULESH~\cite{lulesh}, miniFE~\cite{minife}) maintain constant-size arrays throughout execution, while \emph{variable-size} variants (LAMMPS~\cite{lammps}, WarpX~\cite{warpx}, Nyx~\cite{nyx}) involve particle migration, adaptive mesh refinement, or dynamic load balancing that changes the checkpoint payload between steps.

\paragraph{Task-parallel / DAG.}
These applications decompose computation into tasks with data dependencies. Checkpoint placement must respect task boundaries and data ownership. Chameleon~\cite{chameleon} and DPLASMA~\cite{dplasma} use runtime systems (StarPU~\cite{starpu}, PaRSEC~\cite{parsec}) with algorithm-based fault tolerance (ABFT)~\cite{abft}, while CLAMR~\cite{clamr} and deal.II~\cite{dealii} use adaptive refinement that changes mesh topology.

\paragraph{Pipeline / Staged.}
These applications execute a sequence of distinct computational phases (e.g., mesh generation, solve, post-process). Checkpoints may be placed between stages or within long-running stages. Palabos~\cite{palabos}, Nektar++~\cite{nektar}, and SU2~\cite{su2} have fixed pipelines; RAxML-NG~\cite{raxml-ng} and ExaML~\cite{examl} iterate over a variable-length search space.

\paragraph{Irregular / Adaptive.}
These applications use data-dependent control flow (graph traversal, multigrid cycles, community detection). The checkpoint payload depends on the current refinement level or graph partition. AMG~\cite{amg} and miniVite~\cite{minivite} have fixed graph sizes; SAMRAI~\cite{samrai}, Uintah~\cite{uintah}, and AMReX~\cite{amrex} dynamically refine their meshes.

\paragraph{Application Details.}

Table~\ref{tab:apps} summarizes each application's characteristics.

\noindent
\begin{minipage}{\columnwidth}
\captionof{table}{Benchmark application details. LOC~=~approximate lines of code. Ckpt~Lib~=~checkpoint library used in the reference implementation.}
\label{tab:apps}
\footnotesize
\centering
\setlength{\tabcolsep}{3pt}
\begin{tabular}{@{}llrl@{}}
\toprule
\textbf{Application} & \textbf{Lang.} & \textbf{LOC} & \textbf{Ckpt Lib} \\
\midrule
CoMD~\cite{comd}          & C   &   2K & POSIX \\
LULESH~\cite{lulesh}      & C++ &   5K & SCR~\cite{scr} \\
miniFE~\cite{minife}      & C++ &   3K & FTI~\cite{fti} \\
LAMMPS~\cite{lammps}      & C++ & 300K & Native \\
WarpX~\cite{warpx}        & C++ & 150K & Native \\
Nyx~\cite{nyx}            & C++ &  15K & Native \\
Chameleon~\cite{chameleon} & C++ &   5K & Native \\
DPLASMA~\cite{dplasma}    & C++ &   5K & Native \\
CLAMR~\cite{clamr}        & C++ &  15K & Native \\
deal.II~\cite{dealii}     & C   &   2K & Native \\
Palabos~\cite{palabos}    & C++ &  50K & Native \\
Nektar++~\cite{nektar}    & C   &   3K & Native \\
SU2~\cite{su2}            & C++ & 400K & Native \\
RAxML-NG~\cite{raxml-ng}  & C++ &   5K & Native \\
ExaML~\cite{examl}        & C   &   2K & Native \\
AMG~\cite{amg}            & C   &  30K & FTI~\cite{fti} \\
miniVite~\cite{minivite}  & C++ &   3K & VeloC~\cite{veloc} \\
SAMRAI~\cite{samrai}      & C++ & 200K & Native \\
Uintah~\cite{uintah}      & C   &   3K & Native \\
AMReX~\cite{amrex}        & C++ & 150K & Native \\
\bottomrule
\end{tabular}
\end{minipage}

\paragraph{Selection Criteria.}

Applications were selected to ensure:

\begin{itemize}
    \item \textbf{Coverage}: Every cell in the $4 \times 2$ classification matrix is populated with 2--3 applications.
    \item \textbf{Diversity of checkpoint mechanisms}: The test-bed includes POSIX file I/O, multi-level checkpoint libraries (SCR~\cite{scr}, FTI~\cite{fti}, VeloC~\cite{veloc}), algorithm-based fault tolerance (ABFT)~\cite{abft}, and application-native restart mechanisms.
    \item \textbf{Range of complexity}: From proxy apps (${\sim}$2K~LOC) to production frameworks (${\sim}$400K~LOC).
    \item \textbf{Buildability}: All 20 applications compile and run correctly on the evaluation platform (ARM64 Ubuntu~24.04 with OpenMPI~4.1.6~\cite{openmpi}).
\end{itemize}

%% ---------------------------------------------------------------------------
\subsubsection{Evaluation Approaches}
\label{sec:approaches}

We compare four approaches to adding checkpoint/restart resilience.

\paragraph{Baseline: Unassisted LLM.}

An LLM agent (OpenCode~\cite{opencode} with various backend models) is given the application source code and a prompt requesting VeloC~\cite{veloc} checkpoint integration. The agent modifies the code iteratively: after each modification, automated validation checks correctness. On failure, error logs are fed back to the agent for another attempt. No domain-specific tools or analysis are provided.

We evaluate multiple LLM backends: Claude Sonnet~4~\cite{claude}, Claude Opus~4~\cite{claude}, GPT-4o~\cite{gpt4o}, and Gemini~2.5~Pro~\cite{gemini}, accessed through OpenCode's model-switching capability.

\paragraph{Human-Written Reference.}

For each application, we provide a hand-written checkpointed version that serves as the ground truth. These reference implementations use the checkpoint library specified in the application's configuration (Table~\ref{tab:apps}) and have been validated to correctly recover from process failures. The reference is used to:

\begin{enumerate}
    \item Establish that correct checkpoint/restart \emph{is possible} for each application.
    \item Provide golden output for correctness comparison.
    \item Measure the overhead of checkpointing (runtime, storage).
\end{enumerate}

\paragraph{Transparent Checkpointing (DMTCP/MANA).}

DMTCP~\cite{dmtcp} with the MANA~\cite{mana} MPI plugin provides process-level checkpoint/restart without source modification. We evaluate this approach as a ``zero-effort'' baseline that captures the entire process state (memory, file descriptors, MPI state). While transparent, this approach typically incurs higher overhead and larger checkpoint sizes than application-level approaches~\cite{cr-survey}.

\paragraph{Guard-Agent (Our Approach).}

Guard-Agent augments the LLM agent with domain-specific MCP (Model Context Protocol)~\cite{mcp} tools that provide:

\begin{itemize}
    \item \textbf{Code inspection}: Static analysis of the application's data structures, loop nests, MPI communication patterns, and I/O points.
    \item \textbf{Checkpoint plan generation}: Identification of (a)~which variables to protect, (b)~optimal checkpoint placement in the control flow, and (c)~VeloC API call sequences.
    \item \textbf{Knowledge base}: A curated guide to VeloC~\cite{veloc} API usage, common patterns, and pitfalls.
    \item \textbf{Validation execution}: Automated build, run, failure injection, and output comparison.
    \item \textbf{Failure analysis}: Structured diagnosis of build errors, runtime crashes, and output mismatches.
\end{itemize}

The agent uses these tools within the same iterative loop as the baseline, but with access to structured domain knowledge rather than relying solely on LLM reasoning.

%% ---------------------------------------------------------------------------
\subsubsection{Metrics}
\label{sec:metrics}

\paragraph{Correctness Metrics.}

\begin{itemize}
    \item \textbf{Build success}: Does the modified code compile without errors?
    \item \textbf{Runtime success}: Does the modified code execute and produce output?
    \item \textbf{Recovery success}: After failure injection (process kill after checkpoint), does the application restart and complete?
    \item \textbf{Output correctness}: Does the recovered output match the golden output within the application's tolerance? We employ four comparison methods depending on the application:
    \begin{itemize}
        \item \emph{Numeric}: Element-wise floating-point comparison with configurable relative tolerance (for physics simulations).
        \item \emph{Text}: Line-by-line exact match after filtering non-deterministic lines such as timestamps and timing data.
        \item \emph{Hash}: SHA-256 byte-identical comparison (for binary outputs).
        \item \emph{SSIM}: Structural Similarity Index~\cite{ssim} on HDF5~\cite{hdf5} datasets (for scientific field data).
    \end{itemize}
\end{itemize}

\paragraph{Efficiency Metrics.}

\begin{itemize}
    \item \textbf{Iterations to success} ($N_\text{iter}$): Number of LLM interaction rounds needed to produce a correct resilient version (lower is better).
    \item \textbf{Token consumption} ($T_\text{total}$): Total LLM input + output tokens across all iterations (lower is better).
    \item \textbf{Wall-clock time} ($t_\text{wall}$): Total elapsed time from first prompt to successful validation (lower is better).
    \item \textbf{Success rate} ($S$): Fraction of applications successfully made resilient within the iteration budget (higher is better).
\end{itemize}

\paragraph{Quality Metrics.}

\begin{itemize}
    \item \textbf{Checkpoint data identification}: Precision and recall of data structures identified for protection, compared to the human-written reference.
    \item \textbf{Checkpoint placement accuracy}: Whether checkpoints are placed at correct control flow points (e.g., main loop boundary, after MPI barriers).
    \item \textbf{Restart logic completeness}: Whether the restart path correctly restores all protected state and resumes computation from the checkpoint.
    \item \textbf{Code diff size} ($\Delta_\text{LOC}$): Lines of code added/modified, as a measure of implementation minimality.
\end{itemize}

%% ---------------------------------------------------------------------------
\subsubsection{Experimental Setup}
\label{sec:setup}

\paragraph{Platform.}

All experiments are conducted on an ARM64 (aarch64) development server with 4~cores and 32~GB RAM, running Ubuntu~24.04~LTS. The software stack includes OpenMPI~4.1.6~\cite{openmpi}, GCC~13.3.0, gfortran~13.3.0, CMake~3.28.3, Meson~1.10.2, and Ninja~1.13.0. Checkpoint libraries---VeloC~\cite{veloc} (latest), FTI~1.6~\cite{fti}, SCR~3.1.0~\cite{scr}, and jemalloc~5.3.0~\cite{jemalloc}---are installed under \texttt{\$HOME/.local}.

\paragraph{Validation Pipeline.}

The validation pipeline operates in six stages for each application:

\begin{enumerate}
    \item \textbf{Vanilla build}: Compile the original unmodified source code.
    \item \textbf{Golden run}: Execute the vanilla application without failures to capture reference output.
    \item \textbf{No-recovery verification}: Kill the vanilla application mid-execution and restart it. Verify that it does \emph{not} produce correct output, confirming that resilience requires explicit checkpointing.
    \item \textbf{Checkpointed build}: Compile the checkpointed version (either human reference or LLM-generated).
    \item \textbf{Recovery verification}: Kill the checkpointed application mid-execution and restart it. Verify that it \emph{does} produce correct output, confirming successful recovery.
    \item \textbf{Output comparison}: Compare the recovered output against the golden reference using the application-specific comparison method and tolerance.
\end{enumerate}

\paragraph{Iterative Evaluation Protocol.}

For each application and each approach (baseline, Guard-Agent):

\begin{enumerate}
    \item Start with a clean copy of the vanilla source code.
    \item Present the LLM with the application's checkpoint prompt.
    \item Allow the LLM to modify source files.
    \item Run the validation pipeline.
    \item If validation passes: record success metrics and stop.
    \item If validation fails: extract error logs (compiler output, runtime errors, output diff) and present them to the LLM as feedback for the next iteration.
    \item Repeat until success or the iteration budget ($N_\text{max} = 20$) is exhausted.
\end{enumerate}

\paragraph{Resumption safety.}
If the evaluation is interrupted (e.g., by system restart or agent timeout), the framework re-copies the vanilla source for any incomplete application before resuming, ensuring that partially-modified code from a previous run does not contaminate the next attempt.

\paragraph{Evaluation Dimensions.}

\paragraph{(i) Data Structure Detection.}
We evaluate how well each approach identifies the critical data structures that must be checkpointed. For each application, the human-written reference specifies which variables are protected (documented per-application). We measure:

\begin{itemize}
    \item \textbf{Precision}: What fraction of the variables the LLM chose to checkpoint are actually necessary?
    \item \textbf{Recall}: What fraction of the necessary variables did the LLM correctly identify?
    \item \textbf{False positives}: Variables checkpointed unnecessarily (increases overhead but does not break correctness).
    \item \textbf{False negatives}: Variables missed (causes incorrect recovery---the most critical failure mode).
\end{itemize}

We analyze these metrics across the classification matrix (Table~\ref{tab:classification}) to determine whether certain computation patterns or state characteristics are more challenging for LLM-based analysis.

\paragraph{(ii) Robustness.}
We evaluate robustness along three axes:

\begin{itemize}
    \item \textbf{Data structure complexity}: Fixed-size arrays (CoMD~\cite{comd}) vs.\ dynamically-allocated linked structures (LAMMPS~\cite{lammps} particles) vs.\ framework-managed containers (AMReX~\cite{amrex} \texttt{MultiFab}).
    \item \textbf{Computation flow complexity}: Simple main loops (LULESH~\cite{lulesh}) vs.\ nested iteration with convergence checks (AMG~\cite{amg} V-cycles) vs.\ task DAGs (DPLASMA~\cite{dplasma}) vs.\ multi-phase pipelines (SU2~\cite{su2}).
    \item \textbf{Codebase complexity}: Proxy apps with ${\sim}$2K~LOC and flat file structure vs.\ production frameworks with ${\sim}$400K~LOC, deep directory hierarchies, and complex build systems.
\end{itemize}

For each axis, we report success rate, iteration count, and token usage, and identify failure modes specific to each complexity level.

\paragraph{(iii) Efficiency Improvement.}
We measure how Guard-Agent's domain-specific tools improve efficiency over the unassisted baseline:

\begin{itemize}
    \item \textbf{Iteration reduction}: How many fewer LLM rounds are needed with Guard-Agent?
    \item \textbf{Token reduction}: How much less LLM computation is consumed?
    \item \textbf{Time reduction}: How much faster is the end-to-end process?
    \item \textbf{Failure mode elimination}: Which categories of errors (build failures, incorrect checkpoint placement, missing restart logic) are prevented by Guard-Agent's analysis tools?
\end{itemize}

We further analyze which specific Guard-Agent tools contribute most to efficiency gains: code inspection, checkpoint plan generation, knowledge base queries, or validation feedback.

%% ---------------------------------------------------------------------------
\subsubsection{Reproducibility}
\label{sec:reproducibility}

All evaluation artifacts are included in the repository. Application sources (vanilla and reference checkpointed versions), per-application configuration files (\texttt{app.yaml}), and per-application documentation (computation flow, data structures, checkpoint strategy) are provided for all 20~applications. Dependency installation scripts, evaluation runner scripts, and environment setup are fully automated:

\begin{lstlisting}[language=bash,basicstyle=\small\ttfamily,frame=single,caption={Reproducing the full evaluation on a new system.},label={lst:reproduce}]
# 1. Install system dependencies
sudo ./scripts/install_system_deps.sh

# 2. Install checkpoint libraries (FTI, SCR, jemalloc)
./scripts/install_deps.sh

# 3. Download application source trees
./scripts/install_app_sources.sh

# 4. Set up environment and build infrastructure
source $HOME/.local/env.sh
./setup.sh --clean

# 5. Verify all 20 apps build and validate
./build/run_validate_apps.sh --fresh

# 6. Run full evaluation (baseline + guard-agent)
./build/run_batch.sh --generate-list all > all_apps.txt
./build/run_batch.sh all_apps.txt --mode evaluate \
    --max-iters 20
\end{lstlisting}

%% ---------------------------------------------------------------------------
%% Bibliography entries for the citations used above.
%% Include these in your main .bib file.
%% ---------------------------------------------------------------------------
%%
@misc{comd,
  title = {{CoMD}: Classical Molecular Dynamics Proxy Application},
  author = {{Exascale Co-Design Center for Materials in Extreme Environments}},
  howpublished = {\url{https://github.com/ECP-copa/CoMD}},
  year = {2017}
}

@article{lulesh,
  title = {Hydrodynamics Challenge Problem, {Lawrence Livermore National Laboratory}},
  author = {Karlin, Ian and Broquedis, Jeff and Chamberlain, Bradford L. and others},
  journal = {Technical Report LLNL-TR-490254},
  year = {2013}
}

@inproceedings{minife,
  title = {{MiniFE}: Finite Element Mini-Application},
  author = {Heroux, Michael A. and Doerfler, Douglas W. and Crozier, Paul S. and others},
  booktitle = {Mantevo Project},
  year = {2009},
  howpublished = {\url{https://github.com/Mantevo/miniFE}}
}

@article{lammps,
  title = {Fast Parallel Algorithms for Short-Range Molecular Dynamics},
  author = {Plimpton, Steve},
  journal = {Journal of Computational Physics},
  volume = {117},
  number = {1},
  pages = {1--19},
  year = {1995},
  doi = {10.1006/jcph.1995.1039}
}

@article{warpx,
  title = {{WarpX}: A new exascale computing platform for beam-plasma simulations},
  author = {Vay, Jean-Luc and Almgren, Ann and others},
  journal = {Nuclear Instruments and Methods in Physics Research A},
  volume = {909},
  pages = {476--479},
  year = {2018},
  doi = {10.1016/j.nima.2018.01.035}
}

@article{nyx,
  title = {Nyx: A Massively Parallel {AMR} Code for Computational Cosmology},
  author = {Almgren, Ann S. and Bell, John B. and Lijewski, Mike J. and others},
  journal = {The Astrophysical Journal},
  volume = {765},
  number = {1},
  pages = {39},
  year = {2013},
  doi = {10.1088/0004-637X/765/1/39}
}

@article{chameleon,
  title = {Chameleon: A Hybrid, Proactive Auto-Scaling Mechanism on a Level-Based Queue of Dense Linear Algebra Computations},
  author = {Agullo, Emmanuel and Bramas, B{\'e}renger and Coulaud, Olivier and others},
  journal = {Concurrency and Computation: Practice and Experience},
  year = {2016},
  doi = {10.1002/cpe.3723}
}

@inproceedings{dplasma,
  title = {{DPLASMA}: A New Efficiency in Dense Linear Algebra on Distributed Heterogeneous Hardware Accelerated Systems},
  author = {Bosilca, George and Bouteiller, Aurelien and Danalis, Anthony and others},
  booktitle = {IEEE International Parallel \& Distributed Processing Symposium Workshops},
  year = {2011}
}

@article{clamr,
  title = {{CLAMR}: Cell-based Adaptive Mesh Refinement Mini-App},
  author = {Robey, Robert W. and others},
  howpublished = {\url{https://github.com/lanl/CLAMR}},
  year = {2013}
}

@article{dealii,
  title = {The deal.{II} Library, Version 9.5},
  author = {Arndt, Daniel and Bangerth, Wolfgang and Davydov, Denis and others},
  journal = {Journal of Numerical Mathematics},
  volume = {31},
  number = {3},
  pages = {171--186},
  year = {2023},
  doi = {10.1515/jnma-2023-0089}
}

@article{palabos,
  title = {Palabos: Parallel Lattice {Boltzmann} Solver},
  author = {Latt, Jonas and Malaspinas, Orestis and Kontaxakis, Dimitrios and others},
  journal = {Computers \& Mathematics with Applications},
  volume = {81},
  pages = {334--350},
  year = {2021},
  doi = {10.1016/j.camwa.2020.03.022}
}

@article{nektar,
  title = {Nektar++: An Open-Source Spectral/hp Element Framework},
  author = {Cantwell, Chris D. and Moxey, David and Comerford, Andrew and others},
  journal = {Computer Physics Communications},
  volume = {192},
  pages = {205--219},
  year = {2015},
  doi = {10.1016/j.cpc.2015.02.008}
}

@article{su2,
  title = {{SU2}: An Open-Source Suite for Multiphysics Simulation and Design},
  author = {Economon, Thomas D. and Palacios, Francisco and Copeland, Sean R. and others},
  journal = {AIAA Journal},
  volume = {54},
  number = {3},
  pages = {828--846},
  year = {2016},
  doi = {10.2514/1.J053813}
}

@article{raxml-ng,
  title = {{RAxML-NG}: A fast, scalable and user-friendly tool for maximum likelihood phylogenetic inference},
  author = {Kozlov, Alexey M. and Darriba, Diego and Flouri, Tom{\'a}{\v{s}} and others},
  journal = {Bioinformatics},
  volume = {35},
  number = {21},
  pages = {4453--4455},
  year = {2019},
  doi = {10.1093/bioinformatics/btz305}
}

@article{examl,
  title = {{ExaML} version 3: A tool for phylogenomic analyses on supercomputers},
  author = {Kozlov, Alexey M. and Aberer, Andre J. and Stamatakis, Alexandros},
  journal = {Bioinformatics},
  volume = {31},
  number = {15},
  pages = {2577--2579},
  year = {2015},
  doi = {10.1093/bioinformatics/btv184}
}

@misc{amg,
  title = {{AMG}: Algebraic Multigrid Benchmark},
  author = {{Lawrence Livermore National Laboratory}},
  howpublished = {\url{https://github.com/LLNL/AMG}},
  year = {2017}
}

@inproceedings{minivite,
  title = {Fast and Scalable Graph Community Detection Using {Louvain} Method},
  author = {Ghosh, Sayan and others},
  booktitle = {IEEE International Conference on High Performance Computing, Data, and Analytics},
  year = {2018}
}

@article{samrai,
  title = {Managing Complex Data and Geometry in Parallel Structured {AMR} Applications},
  author = {Hornung, Richard D. and Kohn, Scott R.},
  journal = {Engineering with Computers},
  volume = {22},
  number = {3--4},
  pages = {181--195},
  year = {2006},
  doi = {10.1007/s00366-006-0038-6}
}

@article{uintah,
  title = {Uintah: A Massively Parallel Problem Solving Environment},
  author = {Berzins, Martin and Luitjens, Justin and Meng, Qingyu and others},
  journal = {IEEE International Symposium on High Performance Distributed Computing},
  year = {2010},
  doi = {10.1109/HPDC.2010.5547150}
}

@article{amrex,
  title = {{AMReX}: A Framework for Block-Structured Adaptive Mesh Refinement},
  author = {Zhang, Weiqun and Almgren, Ann and Beckner, Vince and others},
  journal = {Journal of Open Source Software},
  volume = {4},
  number = {37},
  pages = {1370},
  year = {2019},
  doi = {10.21105/joss.01370}
}

@inproceedings{veloc,
  title = {{VeloC}: Very Low Overhead Checkpointing System},
  author = {Nicolae, Bogdan and others},
  booktitle = {IEEE International Conference on Cluster Computing},
  year = {2019},
  doi = {10.1109/CLUSTER.2019.8890984}
}

@inproceedings{scr,
  title = {Scalable Checkpoint/Restart for {MPI}},
  author = {Moody, Adam and Bronevetsky, Greg and Mohror, Kathryn and de Supinski, Bronis R.},
  booktitle = {ACM/IEEE International Conference for High Performance Computing, Networking, Storage and Analysis (SC)},
  year = {2010},
  doi = {10.1109/SC.2010.25}
}

@article{fti,
  title = {{FTI}: High Performance Fault Tolerance Interface for Hybrid Systems},
  author = {Bautista-Gomez, Leonardo and Tsuboi, Seiji and Komatitsch, Dimitri and others},
  journal = {ACM/IEEE International Conference for High Performance Computing, Networking, Storage and Analysis (SC)},
  year = {2011},
  doi = {10.1145/2063384.2063427}
}

@inproceedings{dmtcp,
  title = {{DMTCP}: Transparent Checkpointing for Cluster Computations and the Desktop},
  author = {Ansel, Jason and Arya, Kapil and Cooperman, Gene},
  booktitle = {IEEE International Symposium on Parallel and Distributed Processing (IPDPS)},
  year = {2009},
  doi = {10.1109/IPDPS.2009.5161063}
}

@inproceedings{mana,
  title = {{MANA} for {MPI}: {MPI}-Agnostic Network-Agnostic Transparent Checkpointing},
  author = {Garg, Rohan and Mohan, Apoorve and Sullivan, Michael and Cooperman, Gene},
  booktitle = {ACM International Symposium on High-Performance Parallel and Distributed Computing (HPDC)},
  year = {2019},
  doi = {10.1145/3307681.3325962}
}

@misc{mcp,
  title = {Model Context Protocol},
  author = {{Anthropic}},
  howpublished = {\url{https://modelcontextprotocol.io}},
  year = {2024}
}

@misc{opencode,
  title = {{OpenCode}: Open-source AI coding agent},
  howpublished = {\url{https://github.com/anthropics/opencode}},
  year = {2024}
}

@misc{claude,
  title = {Claude: {AI} Assistant by {Anthropic}},
  author = {{Anthropic}},
  howpublished = {\url{https://www.anthropic.com/claude}},
  year = {2024}
}

@misc{gpt4o,
  title = {{GPT-4o}},
  author = {{OpenAI}},
  howpublished = {\url{https://openai.com/gpt-4o}},
  year = {2024}
}

@misc{gemini,
  title = {{Gemini}: A Family of Highly Capable Multimodal Models},
  author = {{Google DeepMind}},
  year = {2024},
  howpublished = {\url{https://deepmind.google/technologies/gemini/}}
}

@inproceedings{starpu,
  title = {{StarPU}: A Unified Platform for Task Scheduling on Heterogeneous Multicore Architectures},
  author = {Augonnet, C{\'e}dric and Thibault, Samuel and Namyst, Raymond and Wacrenier, Pierre-Andr{\'e}},
  booktitle = {Euro-Par},
  year = {2009}
}

@inproceedings{parsec,
  title = {{PaRSEC}: Exploiting Heterogeneity to Enhance Scalability},
  author = {Bosilca, George and Bouteiller, Aurelien and Danalis, Anthony and others},
  booktitle = {Computing in Science \& Engineering},
  volume = {15},
  number = {6},
  pages = {36--45},
  year = {2013}
}

@article{abft,
  title = {Algorithm-Based Fault Tolerance for Matrix Operations},
  author = {Huang, Kuang-Hua and Abraham, Jacob A.},
  journal = {IEEE Transactions on Computers},
  volume = {C-33},
  number = {6},
  pages = {518--528},
  year = {1984},
  doi = {10.1109/TC.1984.1676475}
}

@article{openmpi,
  title = {Open {MPI}: Goals, Concept, and Design of a Next Generation {MPI} Implementation},
  author = {Gabriel, Edgar and Fagg, Graham E. and Bosilca, George and others},
  journal = {European PVM/MPI Users' Group Meeting},
  pages = {97--104},
  year = {2004}
}

@article{ssim,
  title = {Image Quality Assessment: From Error Visibility to Structural Similarity},
  author = {Wang, Zhou and Bovik, Alan C. and Sheikh, Hamid R. and Simoncelli, Eero P.},
  journal = {IEEE Transactions on Image Processing},
  volume = {13},
  number = {4},
  pages = {600--612},
  year = {2004},
  doi = {10.1109/TIP.2003.819861}
}

@misc{hdf5,
  title = {{HDF5}: Hierarchical Data Format, version 5},
  author = {{The HDF Group}},
  howpublished = {\url{https://www.hdfgroup.org/solutions/hdf5/}},
  year = {2024}
}

@misc{jemalloc,
  title = {jemalloc: A General Purpose malloc(3) Implementation},
  author = {Evans, Jason},
  howpublished = {\url{https://github.com/jemalloc/jemalloc}},
  year = {2024}
}

@article{cr-survey,
  title = {Checkpoint/Restart Approaches for a Thread-Based {MPI} Runtime},
  author = {Hursey, Joshua and Squyres, Jeffrey M. and Mattox, Timothy I. and Lumsdaine, Andrew},
  journal = {Future Generation Computer Systems},
  volume = {25},
  number = {2},
  pages = {206--212},
  year = {2009}
}
